\chapter{Technical Design}
\label{ch:technical-design}

% TODO: expand this chapter to cover the overall architecture of the
% chosen solution -- mechanical, electrical/sensing, and control-software
% architecture -- and, more importantly, why this architecture was chosen.

\section{Conceptual Design Sketch (Phase 1)}
\label{sec:design-sketch-phase1}

\Cref{fig:sentinel_v1_actuant} shows an early-stage hand sketch produced
during Phase 1 of the mechanical design, used to visualise which
functions the rope-pinching mechanism needs to provide before committing
to a detailed CAD model. The sketch identifies the following elements:

\begin{itemize}
    \item An \textbf{encoder}, used to measure the rope's acceleration as
        it moves through the mechanism.
    \item A \textbf{load pin}, integrated into the pivot axle of the
        pinching lever, used to measure the tension on the climbing side
        of the rope (i.e.\ the side leading up to the climber). This
        tension reading provides the information needed to determine
        when the climber has regained a stable position on the wall, so
        that the pinch can be safely released (see
        \cref{sec:fr-post-fall}).
    \item A \textbf{spring} built into the pinching mechanism, providing
        the passive fail-safe behaviour described in
        \cref{sec:fr-failsafe}: the mechanism defaults to a locked/pinched
        state on its own, without requiring power.
    \item A \textbf{servo}, used to actuate the pinching lever -- i.e.\
        to open and close the brake under motor control.
    \item \textbf{Axles for the outer cover}, and a \textbf{hole for the
        carabiner} used to attach the device to the anchor point.
    \item An \textbf{accommodation volume for the electronics}, i.e.\ the
        internal space reserved for the control board and wiring.
    \item A \textbf{potentiometer}, mounted in the pinching mechanism, used
        to sense the lever's current angular position (i.e.\ how open or
        closed the brake currently is).
\end{itemize}

\begin{figure}[H]
    \centering
    \includegraphics[width=0.6\textwidth]{figures/Sentinel v1 actuant.jpg}
    \caption{Phase 1 conceptual sketch of the rope-pinching mechanism,
    used to visualise the device's required functions.}
    \label{fig:sentinel_v1_actuant}
\end{figure}
